\section{Related Work}
\label{sec:related}

This work sits at the intersection of three literatures: hybrid offline-to-online reinforcement learning, sim-to-real dynamics transfer, and data-driven bus holding control.
We survey each in turn and position our contributions against the closest prior work.

\subsection{Hybrid offline-to-online RL and ensemble pessimism}

Hybrid offline-to-online (H2O) RL combines an offline dataset with online interactions to escape the conservative bound of pure offline methods while avoiding the safety cost of pure online learning.
Early works framed the problem as offline RL with a fine-tuning phase: AWAC \citep{nair2020awac} regularises online updates toward the offline behaviour, while RLPD \citep{ball2023efficient} over-samples online transitions to accelerate adaptation.
JSRL \citep{uchendu2023jumpstart} instead uses an offline policy as a guide during online rollouts.
These methods assume a single environment shared between offline and online phases and do not address dynamics gaps.

H2O \citep{niu2022h2o} and its successor H2O+ \citep{liu2024h2oplus} extend the setting to \emph{inaccurate} simulators: the offline dataset may come from the real environment while online rollouts occur in a lower-fidelity simulator.
The central technical device is an importance-sampling (IS) correction on simulator transitions, estimated by a dynamics discriminator, which we review in \Cref{sec:method}.
WSRL \citep{zhou2024wsrl} further proposes a warmup-collect phase and a KL anchor to a policy snapshot to stabilise the offline-to-online transition.
Our method inherits the H2O+ hybrid Bellman structure and optionally uses the WSRL KL anchor, but replaces both the discriminator design and the pessimism mechanism.

Ensemble pessimism is orthogonal to H2O but essential to our target calibration.
EDAC and SAC-$N$ \citep{an2021uncertainty} showed that increasing ensemble disagreement serves as a free pessimism signal, and MSG \citep{ghasemipour2022so} refined this with independent-target variants.
Cal-QL \citep{nakamoto2023calql} re-calibrates conservative Q-learning \citep{kumar2020cql} against a Monte-Carlo reference to prevent over-conservatism during fine-tuning.
These mechanisms were developed in single-environment settings.
Our paper applies an ensemble LCB critic but does not use Cal-QL: instead of Monte-Carlo-based anchoring, we apply a simpler bootstrap one-sided floor (\Cref{eq:qfloor}) tracking the current Q-network's predictions on the offline batch.
We distinguish the two approaches explicitly in \Cref{sec:method:floor}.

\subsection{Sim-to-real transfer via dynamics discrimination}

Two broad paradigms dominate sim-to-real RL: domain randomisation, which trains a policy robust to a family of perturbed dynamics \citep{tobin2017domain, peng2018sim}, and dynamics correction, which actively estimates and corrects the sim/real gap.
Our work belongs to the latter.

GARAT \citep{desai2020garat} and its predecessor grounded action transformation \citep{hanna2017grounded} learn an action transformation to make sim behave like real; these require paired data and do not extend naturally to transit, where paired trajectories are unavailable.
DARC \citep{eysenbach2021off} instead learns a classifier $\disc(s,a,s')$ distinguishing real from sim transitions and uses it to reweight simulator rollouts via IS; H2O and H2O+ both adopt DARC-style discriminators.
\Cref{sec:method} argues that DARC's reliance on action-conditioned transitions fails in sparse-action transit control, and proposes an InfoNCE-based action-invariant alternative.

Contrastive representation learning has been used in RL for auxiliary tasks---CURL \citep{laskin2020curl} for image encoders and CPC++ \citep{mazoure2020deep} for dynamics---but not, to our knowledge, directly as a sim-to-real IS estimator.
Our use of InfoNCE as a likelihood-ratio estimator is grounded in the classical identification of the optimal NCE score with the log-density ratio \citep{oord2018representation, ma2018nce, poole2019variational}---a connection that, to our knowledge, has not been exploited for multi-fidelity sim-to-sim IS correction in RL.

\subsection{Bus holding control with reinforcement learning}

Analytical bus holding control has a long pedigree, beginning with \citet{daganzo2009headway} and \citet{xuan2011dynamic}, who derived closed-form holding strategies to suppress bunching under simplified demand models.
The RL literature has grown in the last five years: \citet{alesiani2018reinforcement} study RL-based holding for high-frequency services, \citet{wang2020real} develop a multi-agent deep RL framework for dynamic holding, and later studies extend the setting to stochastic event-driven holding and multi-line shared-corridor control \citep{liu2023deepholding, wang2023multiobjective}.

All of these works treat their simulator as ground truth: the agent is trained and evaluated in the same environment, and the cross-fidelity gap (training in a cheaper simulator than the evaluation environment) is discussed only anecdotally or as future work.
To our knowledge, this paper is the first to pose bus holding control as an explicit hybrid offline-to-online problem with a dynamics gap, and correspondingly the first to diagnose where off-the-shelf H2O+ components break down on transit-specific MDP structure.
Offline and offline-to-online RL for transit operations remain comparatively underexplored relative to simulator-trained online RL.
Our contribution is therefore not another simulator-only holding controller, but a diagnosis of what happens when historical data, a cheap online simulator, and a higher-fidelity evaluation simulator are intentionally separated.
